%% kapa: exam=1 round=25 year=2014 period=2 form=A subject=accounting subject_ko=회계학 seq=16 number=56
%% source: Q-Net 감정평가사 기출문제 (cst00309) artlSeq=5101207 file=extracted/1008/5101207/2014년 제25회 감정평가사 1차시험 기출문제/2교시(A형).hwp
%% conversion: kordoc 4.24.51 -> tools/gen1.py (auto); review=proofread-sonnet; flags=none
\begin{question}{56}
수익에 관한 설명으로 옳지 않은 것은?
\begin{choices}
  \item 인도된 재화의 결함에 대하여 정상적인 품질보증범위를 초과하여 책임을 지는 경우에는 수익을 인식하지 않는다.
  \item 판매대금의 회수가 구매자의 재판매에 의해 결정되는 경우에는 수익을 인식하지 않는다.
  \item 설치조건부 판매에서 계약의 중요한 부분을 차지하는 설치가 아직 완료되지 않은 경우에는 수익을 인식하지 않는다.
  \item 판매자가 판매대금의 회수를 확실히 할 목적만으로 해당 재화의 법적 소유권을 계속 가지고 있더라도 소유에 따른 중요한 위험과 보상이 이전되었다면 해당 거래를 판매로 보아 수익을 인식한다.
  \item 수익은 기업이 받았거나 받을 경제적효익의 총유입을 포함하므로, 판매세, 특정재화나 용역과 관련된 세금, 부가가치세와 같이 제3자를 대신하여 받는 금액도 기업에 유입된 경제적효익이므로 기타수익에 포함시킨다.
\end{choices}
\end{question}
\answer{5}
